\documentclass[11pt]{article}

\usepackage[margin=1in]{geometry}
\usepackage{microtype}
\usepackage{graphicx}
\usepackage{booktabs}
\usepackage{array}
\usepackage{longtable}
\usepackage{xcolor}
\usepackage{listings}
\usepackage{amsmath}
\usepackage{amssymb}
\usepackage[hidelinks,colorlinks=true,linkcolor=blue!50!black,citecolor=blue!50!black,urlcolor=blue!50!black]{hyperref}
\usepackage[capitalize,noabbrev]{cleveref}
\usepackage{natbib}
\usepackage{titlesec}
\usepackage{enumitem}

\definecolor{codebg}{RGB}{245,245,245}
\lstset{
  basicstyle=\ttfamily\small,
  backgroundcolor=\color{codebg},
  columns=fullflexible,
  breaklines=true,
  breakindent=0pt,
  frame=single,
  framesep=4pt,
  rulecolor=\color{gray!40},
  xleftmargin=4pt,
  xrightmargin=4pt,
}

% This document is unusually dense with long inline \texttt identifiers (config keys,
% tensor names, class names) that TeX's default line breaker cannot hyphenate, which
% produces overfull lines with ordinary tolerance settings. Loosening interword stretch
% is the scalable fix, rather than hand-rewording every paragraph that hits one.
\tolerance=1200
\emergencystretch=3em
\hbadness=10000
\frenchspacing
\setlength{\parskip}{0.4em}
\setlength{\parindent}{0pt}

\titleformat{\section}{\large\bfseries}{\thesection}{0.6em}{}
\titleformat{\subsection}{\bfseries}{\thesubsection}{0.6em}{}
\titleformat{\subsubsection}{\bfseries\itshape}{\thesubsubsection}{0.6em}{}

\title{\bfseries unfuse: Splitting, Verifying, and Quantizing \\ Multimodal Language Models}
\author{Sartaj Bhuvaji}
\date{October 2026}

\begin{document}
\maketitle

\begin{abstract}
\noindent
A multimodal language model is an encoder bolted onto a language model: a vision tower
that turns pixels into vectors, feeding a decoder that was, underneath, already a text
model. \texttt{unfuse} takes that apart --- keep the language half, keep the encoder half,
or keep both separately --- as a tensor filter and a rename, with no retraining and no
quality loss claimed without a measurement to back it. This report documents three
case studies built with it (GLM-4.6V-Flash, Qwen3.8-27B, and a requantized Gemma~4 31B),
the verification methodology that makes ``no quality loss'' a checked claim rather than
an assumption, and what each investigation got wrong before it got right. The headline
result is incidental to the splitting work: quantizing a model for serving measurably
degrades its speculative-decoding acceptance rate, independent of wall-clock, monotonic
across three points. This is a technical report, not a claim of novel method --- the value
is in what was measured and what broke along the way.
\end{abstract}

\section{Introduction}
\label{sec:intro}

A model that reads images is two things glued together: an encoder that turns pixels
into vectors, and an ordinary language model that consumes them. The two halves are
trained together, shipped as one checkpoint, and in many deployments only one of them
is ever used. A text-only fine-tuning run still pays to load a vision tower it never
calls. A \texttt{llama.cpp} server that only wants the language half still carries the
encoder's weights in every GGUF quant. \texttt{unfuse} takes the two apart: keep the
language model, keep the encoder, or keep both separately. It is a tensor filter and a
rename, with no retraining and no distillation.

The split itself turns out to be the easy part. Most multimodal architectures separate
cleanly along existing module boundaries --- the hard part, and the actual subject of
this report, is \emph{proving} the result is unchanged rather than assuming it from the
architecture diagram. Section~\ref{sec:methodology} describes the two verification bars
this project settled on, and why a single tolerance threshold is the wrong tool for
either half of the job. Sections~\ref{sec:glm}--\ref{sec:gemma} are three case studies,
in the order they were done, each building on a lesson the previous one left behind:
GLM-4.6V-Flash (Section~\ref{sec:glm}) establishes that the split can be bit-exact, not
just close, and that most of a widely-quoted memory win was never about the weights at
all. Qwen3.8-27B (Section~\ref{sec:qwen}) is a much larger, much more popular model
whose vision tower is a rounding error --- 1.658\% of the checkpoint --- so the
motivation for splitting it has nothing to do with size. Gemma~4 31B
(Section~\ref{sec:gemma}) is the model that broke the premise of this whole project: its
encoder is a single matmul, not a tower, so there is almost nothing to remove, and the
investigation pivots from splitting to quantization. That pivot produced this report's
one genuinely novel measurement, not a splitting result at all: quantizing a model for
serving measurably costs speculative-decoding throughput, independent of wall-clock,
scaling monotonically with how aggressively the target is quantized.

Section~\ref{sec:lessons} collects the lessons that generalize past any one model ---
config validators, serving-engine quirks, and infrastructure traps that cost real GPU
time to find. None of it is dressed up as more than it is. This is a technical report of
what was built, measured, and broken, not a claim that any of it is a novel method.

\section{Methodology: Two Verification Bars}
\label{sec:methodology}

Two different operations run through \texttt{unfuse}, and they cannot share a
correctness bar.

\textbf{A split is pure plumbing.} Dropping a tensor, renaming a prefix, and
reconstructing a config do no arithmetic on the weights that survive. The result either
reproduces the original model's output bit-for-bit, or the extraction is broken --- there
is no third option, and no tolerance to tune. \texttt{unfuse verify} holds that line: it
compares logits between the original and the extracted model on a fixed prompt set and
requires \texttt{max|d| = 0.000e+00}. A unit test pins the bar there so it cannot drift.

\textbf{Quantization is the opposite case.} It rounds every weight it touches, so a
non-zero difference from the original is not a bug, it is the entire point of the
operation. Scoring it against the same bit-exact bar is meaningless by construction. The
tempting fix --- pick an absolute tolerance and call anything under it a pass --- does
not survive contact with a real pipeline: nobody can independently justify the number,
and the moment a threshold exists, it will eventually be loosened until whatever
checkpoint is on the bench that day passes it. \texttt{unfuse agreement} uses a
comparative bar instead. Whoever published the quantized build being modified already
decided how much divergence from the original is acceptable --- they shipped benchmark
numbers next to it. Score their checkpoint against the original, score this project's
checkpoint the same way, and require the latter to be no worse:

\begin{lstlisting}
unfuse agreement google/gemma-4-31B-it ~/models/mine \
    --baseline nvidia/Gemma-4-31B-IT-NVFP4
\end{lstlisting}

The two modules deliberately share no code. Merging them would break both: the exact bar
would reject every legitimate quantization (all of them move some bits), and the
tolerant bar would wave through a broken extraction that happened to land inside
whatever threshold survived the merge.

A strict \emph{no-worse-than-baseline} gate is the right test only when the change under
evaluation is not supposed to cost anything. It is the wrong test the moment that
assumption breaks --- Section~\ref{sec:gemma} adds FP8 attention on top of NVIDIA's own
NVFP4 baseline, so it cannot pass a strict gate by construction. For exactly that case
\texttt{Agreement.marginal\_cost} reports the \emph{increment} over the baseline instead
of a pass/fail verdict, which is the only honest thing to report when a change is known
in advance to add some cost in exchange for something else.

\textbf{The bar only works if the yardstick is not broken, and checking that costs
real effort.} The first scoring run against the Gemma~4 31B quantization (full account in
Section~\ref{sec:gemma}) produced numbers saying this project's checkpoint \emph{beat}
NVIDIA's own build. It was noise: the BF16 reference itself was emitting
\texttt{France is France is France is\ldots} because \texttt{gemma-4-31B-it} is
instruction-tuned and the scoring script was feeding it raw completion prompts instead of
the chat template. Scoring in that degenerate regime inflates divergence by roughly an
order of magnitude in arbitrary directions, which is exactly how a \emph{more} quantized
checkpoint came out looking closer to the reference. \texttt{agreement} grew a
\texttt{--chat} mode and a fixed prompt set as a direct result; \texttt{verify} kept its
raw prompts, because an extraction has to be bit-exact on any input, degenerate or not ---
that is a different claim from scoring quantization damage, so the two tools diverge
again at exactly the place the exact/comparative split already drew. The general lesson
generalizes past this one incident: a result that happens to favor the work already done
is the one to distrust most, and it is worth the extra step to find out why before
reporting it.

\section{Case Study: GLM-4.6V-Flash}
\label{sec:glm}

\texttt{zai-org/GLM-4.6V-Flash} (\texttt{model\_type: glm4v},
\texttt{Glm4vForConditionalGeneration}) is the model this project's verification
methodology was built against, and the one result published without qualification.

\subsection{What Is Actually There}

The checkpoint's own \texttt{model.safetensors.index.json} reports a \texttt{total\_size}
that, read as bytes, implies roughly 5.1B parameters. It is not bytes --- for this repo
that field reports the parameter count directly. The real figure, confirmed against the
sum of every tensor shape and against the shard files' actual byte sizes, is
10{,}292{,}777{,}472 parameters.

\begin{table}[h]
\centering
\small
\begin{tabular}{lrr}
\toprule
Group & Params & Share \\
\midrule
\texttt{visual.blocks.0--23} (ViT, 1536d, 12 heads) & 679,550,976 & 6.602\% \\
\texttt{visual.merger} & 185,081,856 & 1.798\% \\
\texttt{visual.downsample} & 25,169,920 & 0.245\% \\
\texttt{visual.patch\_embed.proj} & 1,807,872 & 0.018\% \\
\texttt{visual.embeddings.position\_embedding} & 884,736 & 0.009\% \\
\texttt{visual.\{post\_conv,post\}\_layernorm} & 3,072 & 0.000\% \\
\midrule
\textbf{DELETE --- \texttt{model.visual.*}, 181 tensors} & \textbf{892,498,432} & \textbf{8.671\%} \\
\texttt{language\_model.layers.0--39} & 8,158,760,960 & 79.267\% \\
\texttt{language\_model.embed\_tokens} & 620,756,992 & 6.031\% \\
\texttt{lm\_head} (untied) & 620,756,992 & 6.031\% \\
\texttt{language\_model.norm} & 4,096 & 0.000\% \\
\midrule
\textbf{KEEP --- 523 tensors} & \textbf{9,400,279,040} & \textbf{91.329\%} \\
\bottomrule
\end{tabular}
\caption{GLM-4.6V-Flash parameter breakdown, verified against tensor shapes.}
\end{table}

The surgery is five steps:
\begin{enumerate}[nosep,leftmargin=1.5em]
\item drop every key under \texttt{model.visual.} (181 tensors);
\item rename the \texttt{model.language\_model.} prefix to \texttt{model.};
\item keep \texttt{lm\_head.weight} (embeddings are untied);
\item rebuild the config via \texttt{Glm4Config}, dropping \texttt{mrope\_section} and
  \texttt{vision\_config}; and
\item set \texttt{architectures} to \texttt{["Glm4ForCausalLM"]} explicitly ---
  \texttt{Glm4Config} will not populate it, and without it
  \texttt{AutoModelForCausalLM} has nothing to dispatch to.
\end{enumerate}
The vocabulary needs no changes: all 151,552 tokens, including the image/video special
tokens, carry over whether or not anything multimodal is ever sent to the model again.

\subsection{Verification: The RoPE Question}

The one real risk in this split was positional encoding. GLM-4.6V-Flash's text config
carries \texttt{mrope\_section: [8, 12, 12]}, multimodal RoPE that splits the rotary
budget across temporal, height, and width axes --- $8+12+12=32$, which is
$\text{head\_dim } 128 \times \text{partial\_rotary\_factor } 0.5 \div 2$. For text-only
input, all three axes carry an identical position index, so mRoPE should mathematically
collapse to standard RoPE. That is an argument from the config, not a measurement, and
this project does not ship arguments from the config without checking them.

It checks out exactly. Comparing logits between the original multimodal model and the
extracted text-only model:

\begin{lstlisting}
max|d| = 0.000e+00   across six prompts, argmax_match=True on all
max|d| = 0.000e+00   at 1,207 tokens
\end{lstlisting}

Zero divergence, not near-zero divergence, at over a thousand tokens of generation. This
is the result the rest of this project's verification discipline is built on: an
extraction with the right discipline behind it does not approximate the original model,
it reproduces it.

\subsection{The Inverse Operation}

Keeping only the vision tower is the mirror image, and mechanically much simpler. The
class \texttt{Glm4vVisionModel} already exists in \texttt{transformers}, and its
submodules map exactly onto what sits under \texttt{model.visual.}, so stripping the
prefix is the whole surgery --- no architecture retarget, no RoPE question. Two facts
make it runnable on a laptop rather than a rented GPU: all 181 vision tensors live in
shard~4 only (confirmed via the index's \texttt{weight\_map}), so the download is
4.66~GB rather than the full 20.59~GB, and the result is 1.78~GB. Run on an RTX~2060
(6.44~GB, SM75, so float32, no bf16 support) with 7~GiB of system RAM:

\begin{lstlisting}
181 tensors, 892,498,432 params, 1.785 GB
byte-compared 181 tensors against the source shard -> all bit-identical
max|d| = 0.000e+00   (64, 4096)
448x448 image -> 1024 patches -> 256 soft tokens x 4096d
\end{lstlisting}

Verification here is deliberately weaker than the removal direction: the full
multimodal model does not fit on a laptop, so this compares a directly-loaded reference
against the \texttt{from\_pretrained} checkpoint rather than against the full original.
It guards the plumbing --- config reconstruction, key naming, the round trip --- not the
arithmetic, which is identical by construction given the byte-level tensor comparison
above. One WSL-specific note: the run reported 7.32~GB peak VRAM on a 6.44~GB card.
WSL2's WDDM driver lets CUDA spill into shared system memory; a deliberate 7~GB
allocation succeeding confirmed this rather than leaving it as a guess. The same code
would OOM on native Linux with the same card.

The output of this direction is a feature extractor, not a model --- no language head, no
classifier, just embeddings aligned to GLM-4's input space. This asymmetry is worth
naming explicitly: removing vision leaves a complete, working model, because a decoder
with its own \texttt{lm\_head} stands alone. Removing text leaves a component that needs
something else to consume its output.

\subsection{GGUF: mmproj Towers Do Not Take K-Quants}

Built with \texttt{llama.cpp} at master (\texttt{GLM4} architecture support). All five
quantizations load and generate coherently; no perplexity comparison against the bf16
original was run for this direction.

\begin{table}[h]
\centering
\small
\begin{tabular}{lr}
\toprule
Quant & Size (GB) \\
\midrule
F16 & 18.81 \\
Q8\_0 & 10.00 \\
Q6\_K & 8.27 \\
Q5\_K\_M & 7.05 \\
Q4\_K\_M & 6.17 \\
\bottomrule
\end{tabular}
\caption{Text-only GGUF quantization ladder, GLM-4.6-Flash-text.}
\end{table}

The vision half does not take the same ladder. \texttt{convert\_hf\_to\_gguf.py
--mmproj} only accepts \texttt{f32}, \texttt{bf16}, \texttt{f16}, and \texttt{q8\_0} ---
\texttt{llama-quantize} and the K-quant ladder apply only to the text GGUF. This is
confirmed, not assumed, against what publishers actually ship: unsloth's Gemma~3 and
Gemma~4 mmproj builds are BF16/F16/F32 only, and ggml-org's Qwen2.5-VL mmproj ships
Q8\_0/F16 only. Q8\_0 at 1.03~GB is the floor for this tower; there is no Q4\_K\_M
projector, for any model. Building it also needed a staging workaround:
\texttt{--mmproj} requires the full multimodal architecture and cannot read a stripped
\texttt{Glm4vVisionModel} checkpoint directly, so the build stages a directory with the
original \texttt{config.json} plus an index whose \texttt{weight\_map} references only
shard~4 --- no second download. \texttt{llama.cpp} has no separate \texttt{GLM4V}
architecture (only \texttt{CHATGLM}, \texttt{GLM4}, \texttt{GLM4\_MOE}, with
\texttt{glm4v} existing only as a vision-projector type), so both the text-only
extraction and GLM-4.6V-Flash's own text tower convert to the same \texttt{glm4} arch,
and the mmproj should pair with either. This is verified from GGUF metadata, not by
running end-to-end multimodal inference, which needed more RAM than the build machine
had.

\subsection{Most of the Memory Win Is the Profiler, Not the Weights}
\label{sec:glm-memory}

The published model card quoted a memory win for the text-only split without ever
benchmarking the parent it was derived from --- only the derivative had been measured.
Closing that gap, on one A100-SXM4-40GB under vLLM~0.27.1 at
\texttt{--max-model-len 32768}, the split serves 479,280 KV tokens against the parent's
310,512: +54\%. The vision tower is only 1.79~GB of 20.59, so that gap cannot be
weights. vLLM's own memory profile says where it actually goes:

\begin{table}[h]
\centering
\small
\begin{tabular}{lrrr}
\toprule
 & Parent & Parent, \texttt{video: 0} & Split \\
\midrule
Weights + non-torch & 19.67 GiB & 19.67 GiB & 17.85 GiB \\
Peak activation & 4.82 GiB & 0.97 GiB & 0.21 GiB \\
KV cache & 11.85 GiB & 15.70 GiB & 18.28 GiB \\
KV tokens & 310,512 & 411,552 & 479,280 \\
Engine init & 116 s & 23 s & 11 s \\
\bottomrule
\end{tabular}
\caption{vLLM memory profile, GLM-4.6V-Flash parent vs.\ text-only split.}
\end{table}

The parent's 4.82~GiB peak activation figure is vLLM profiling the multimodal input
path against one maximum-size \emph{video} input --- not an image --- and sizing a
66,242-token encoder cache against it. That memory is reserved for the life of the
server whether or not a multimodal request is ever sent, and it costs 105~s of the
parent's startup time as well. Which means 60\% of the headline gain
(101,040 of 168,768 tokens) is available on the \emph{parent}, unmodified, by passing
\texttt{--limit-mm-per-prompt '\{"video": 0\}'} at launch, and no surgery at all. The
split is worth +16\% over a parent configured that way, plus 1.73~GiB of weights that no
launch flag reaches. A card claiming the full +54\% against a default-configured parent
is not technically wrong, but it is comparing against a misconfiguration rather than the
best achievable baseline, and both numbers belong in an honest card. (Do not pass
\texttt{\{"video": 0, "image": 0\}} together --- vLLM~0.27.1 then skips loading the
tower entirely and dies in profiling with an \texttt{AttributeError}.)

Both checkpoints store KV at an identical 40,960 bytes/token (40 layers, 2 KV heads,
head\_dim 128, bf16), which is the arithmetic check behind every token figure above ---
tokens must equal GiB $\times$ 26{,}214.4 --- and confirms none of the measured gain is a
cheaper cache; it is entirely reserved-memory accounting.

\subsection{Throughput Did Not Move, and That Is the Result}

Eight benchmark rows (two checkpoints, two engines, concurrency~1 and~16): every
concurrency-1 row lands between 61.8 and 63.7~tok/s, every concurrency-16 row between
962.7 and 994.9 --- a 3\% spread across the entire table. Both checkpoints read the same
18.80~GB per decoded token, because decode never touches the resident vision tower
regardless of whether it is present, so both sit against the same bandwidth ceiling on
the A100's 1.555~TB/s and both land around 77\% of it. The trap to avoid when computing
that ceiling is charging the parent for its whole 20.59~GB checkpoint rather than the
bytes decode actually reads --- that lowers its ceiling by 9\% and manufactures a
throughput advantage for the split that the measurements do not actually show.

Two comparability checks were necessary before any of the above could be reported
honestly. vLLM gave both models \texttt{max\_num\_seqs=256} (clamped from
\texttt{max\_num\_batched\_tokens=2048}), but only prints the clamp at \texttt{DEBUG}
log level, so confirming the two runs were actually comparable took a second start with
identical flags --- two unlogged values are unverified, not equal, until checked.
SGLang's own derived values were 3,085 for the parent and 4,096 for the split; unequal,
but both roughly 200$\times$ the tested concurrency, so nothing was throttled in either
run. The Qwen~3.8-27B benchmark in Section~\ref{sec:qwen} was not so lucky. And warmup
earned its keep on exactly one of the eight rows: SGLang on the split read 50.7~tok/s
cold against 63.3~tok/s warm (+24.9\%, with TPOT identical at 15.58~ms). Compared
against SGLang's warmed parent figure of 63.7, the cold number would have published the
split as 20\% \emph{slower} than the model it was derived from.

\section{Case Study: Qwen3.8-27B}
\label{sec:qwen}

\texttt{Qwen/Qwen3.8-27B} (\texttt{model\_type: qwen3\_5},
\texttt{Qwen3\_5ForConditionalGeneration}) is the opposite case from
Section~\ref{sec:glm} in almost every dimension that matters. It is much larger, its
vision tower is a rounding error of the checkpoint rather than a meaningful fraction of
it, and the reason to split it has nothing to do with size at all: at the time of
writing it had 2.95M downloads/month and 12,691 likes on the Hub, roughly 15$\times$ the
demand of GLM-4.6V-Flash. The split is worth doing because the model is popular, not
because it is heavy.

\subsection{What Is Actually There}

The whole checkpoint was measured by reading all 18 shard headers over HTTP range
requests --- about 200~KB fetched, no weights downloaded, before deciding whether the
split was worth attempting at all.

\begin{table}[h]
\centering
\small
\begin{tabular}{lrrr}
\toprule
Group & Params & Share & GB \\
\midrule
\texttt{model.language\_model.*} (64 layers) & 25,624,600,064 & 92.236\% & 51.25 \\
\texttt{lm\_head} (untied) & 1,271,398,400 & 4.576\% & 2.54 \\
\texttt{mtp.*} (fused MTP head, 15 tensors) & $\sim$424,694,784 & $\sim$1.529\% & 0.85 \\
\midrule
\textbf{DELETE --- \texttt{model.visual.*}, 333 tensors} & \textbf{460,730,096} & \textbf{1.658\%} & \textbf{0.92} \\
\midrule
\textbf{Total} & \textbf{27,781,427,952} & 100\% & \textbf{55.56} \\
\bottomrule
\end{tabular}
\caption{Qwen3.8-27B parameter breakdown, measured from shard headers only.}
\end{table}

The tower is 1.658\% of the checkpoint, against GLM-4.6V-Flash's 8.671\%. Removing it
saves 0.92~GB out of 55.56 --- this is the same lesson Gemma~4 teaches more sharply in
Section~\ref{sec:gemma}, and the reason to run \texttt{unfuse inspect} before assuming
any model is worth splitting at all. The actual win here is packaging: the result loads
as a plain causal LM, with no vision dependencies and no image-preprocessing code path,
not a smaller file.

\subsection{A False Alarm: \texttt{deepstack\_visual\_indexes}}

The config field \texttt{vision\_config.deepstack\_visual\_indexes} reads like it should
complicate the split --- it names tower layers whose features get injected into the
language model at several depths, which sounds like vision reaching into the LLM rather
than staying cleanly separable. It does not turn out that way.
\textbf{Zero tensors in the checkpoint carry a deepstack name.} The only trees present
are \texttt{model.visual.*}, \texttt{model.language\_model.*}, \texttt{lm\_head}, and
\texttt{mtp.*} --- all 1,199 tensors accounted for across those four, nothing left
unclassified. The indices the config names are resolved entirely inside the visual
tree, and the split stays clean. Worth recording precisely because the config
\emph{looks} like a problem on first read, and the answer cost exactly one probe to
settle rather than an afternoon of architecture study.

\subsection{A Hybrid, Not a Transformer}

\texttt{layer\_types} is 48 \texttt{linear\_attention} layers plus 16
\texttt{full\_attention} layers, with \texttt{full\_attention\_interval: 4} --- every
fourth layer is full attention, the rest are linear, alongside
\texttt{linear\_conv\_kernel\_dim: 4}, \texttt{linear\_num\_key\_heads: 16},
\texttt{linear\_num\_value\_heads: 48}, and \texttt{mamba\_ssm\_dtype: float32}. Vanilla-
transformer assumptions about KV cache sizing, GGUF conversion, and engine support do
not transfer to this stack, and several of the traps below are exactly that assumption
failing in a new place.

\subsection{The MTP Head Has No Home in \texttt{transformers}}

\texttt{mtp.*} is a fused multi-token-prediction speculative head: 15 tensors,
$\sim$424.7M parameters, \texttt{mtp\_num\_hidden\_layers: 1}. \texttt{transformers~5.16}
defines \texttt{Qwen3\_5ForCausalLM}, \texttt{Qwen3\_5TextConfig}, and
\texttt{Qwen3\_5VisionModel}, but no MTP class for this family, so
\texttt{Qwen3\_5ForCausalLM} reports these fifteen tensors as unexpected keys and
silently drops them. The split keeps them by default anyway --- dropping them removes a
real capability, and the MTP-bearing community GGUFs of this model are among its most
downloaded builds, so silently shipping a build without it would be a regression users
would not be warned about. \texttt{--drop-mtp} exists as an opt-in for anyone who wants
the clean load instead.

Loading the split reports 26,895,998,464 parameters against 27,320,697,856 actually in
the checkpoint. The difference, 424,699,392, is exactly the MTP head --- confirming by
subtraction that \texttt{transformers} is loading everything else and only silently
ignoring this one component. vLLM~0.27.1 does ship
\texttt{model\_executor/models/qwen3\_5\_mtp.py} and registers \texttt{"Qwen3\_5MTP"}, so
this looked recoverable there. It was not, on either path tested:

\begin{table}[h]
\centering
\small
\begin{tabular}{ll}
\toprule
Config \texttt{model\_type} & Result \\
\midrule
\texttt{qwen3\_5\_text} (ours) & \texttt{NotImplementedError: Unsupported speculative method: 'mtp'} \\
\texttt{qwen3\_5} (forced) & resolves \texttt{Qwen3\_5MTP}, then an \texttt{AssertionError} on a shape mismatch \\
\bottomrule
\end{tabular}
\caption{Both attempted paths to recovering the MTP head under vLLM.}
\end{table}

The cause lives in \texttt{vllm/config/speculative.py:516}: the MTP auto-detect fires
only for \texttt{model\_type in ("qwen3\_5", "qwen3\_5\_moe")}, and that code path then
reads the \emph{multimodal} config layout, expecting a nested \texttt{text\_config}. A
flattened text-only config makes it fall back to a default \texttt{hidden\_size} of 4096
against the real value of 5120. So the head costs 0.85~GB and does nothing today ---
kept anyway, since the alternative is deleting a capability that cannot be recovered
without a fresh 55~GB download, and the only cost of keeping it is the 0.85~GB.

\subsection{The Surgery, and Verifying It Three Times on Two Machines}

The surgery:
\begin{enumerate}[nosep,leftmargin=1.5em]
\item drop every key under \texttt{model.visual.} (333 tensors);
\item rename \texttt{model.language\_model.} to \texttt{model.};
\item keep \texttt{lm\_head.weight} (untied);
\item build the config from \texttt{text\_config} via \texttt{Qwen3\_5TextConfig},
  dropping \texttt{image\_token\_id}, \texttt{video\_token\_id},
  \texttt{vision\_start\_token\_id}, \texttt{vision\_end\_token\_id}, and
  \texttt{language\_model\_only}; and
\item set \texttt{architectures} to \texttt{["Qwen3\_5ForCausalLM"]} explicitly --- the
  same trap \texttt{Qwen3\_5TextConfig} shares with \texttt{Glm4Config} in
  Section~\ref{sec:glm}.
\end{enumerate}
All six non-default architecture keys (\texttt{attn\_output\_gate},
\texttt{full\_attention\_interval}, \texttt{mamba\_ssm\_dtype},
\texttt{mtp\_num\_hidden\_layers}, \texttt{mtp\_use\_dedicated\_embeddings},
\texttt{output\_gate\_type}) were verified to survive the round trip through
\texttt{Qwen3\_5TextConfig} rather than assumed to.

The split ran three separate times --- once locally, twice on a rented H100, each from a
fresh download of the parent --- and all three produced an identical
27,320,697,856~params / 54.64~GB result, all three passing the byte-level check on all
866 kept tensors. Reproduction was checked, not assumed, here as everywhere else in this
project. Verification for this model is tensor hashes, not logits:
\texttt{unfuse.verify.compare\_logits} is the right tool when the surgery could plausibly
change arithmetic, but this one only drops and renames, so a byte equality check over
every kept tensor proves every possible input produces an identical output, where
matching logits only proves it for the specific prompts tried. Byte hashing is also the
only option on a machine that cannot hold a 27B model in RAM at once.

Two implementation details cost real time here. \textbf{numpy has no bfloat16}: calling
\texttt{.numpy()} on a bfloat16 tensor raises \texttt{unsupported ScalarType BFloat16}
on exactly the dtype the checkpoint is stored in, which meant the first version of the
verifier could not check the one thing it existed to check. The fix is to hash
\texttt{t.contiguous().view(torch.uint8).numpy().tobytes()} instead --- since the claim
under test is byte equality, not numeric equality, this also happens to work unmodified
for fp8 checkpoints. And \textbf{the output sharding should not mirror the input
sharding}: reading shard~$N$ and immediately writing shard~$N$ holds a whole $\sim$3~GB
shard in RAM and dies partway through on a small machine, leaving a checkpoint on disk
that still looks plausible but is truncated. Buffering by bytes and flushing on a fixed
budget instead means output sharding does not match input sharding --- which is fine,
because the index file is the authoritative map either way.

\subsection{Decode Is Bandwidth-Bound, and the Numbers Say So}
\label{sec:qwen-bandwidth}

Worth computing before attributing any single-stream throughput result to engine choice
or anything else. An H100~SXM5 has 3.35~TB/s of HBM3 bandwidth, and a dense model at
batch~1 reads every weight once per decoded token:

\begin{table}[h]
\centering
\small
\begin{tabular}{lrrrrr}
\toprule
 & Weights & Floor/token & Ceiling & Measured & Of ceiling \\
\midrule
Qwen3.8-27B (parent) & 55.56 GB & 16.59 ms & 60.3 tok/s & 48.8 tok/s & $\sim$81\% \\
Text-only split & 54.64 GB & 16.31 ms & 61.3 tok/s & 47.2 tok/s & $\sim$77\% \\
\bottomrule
\end{tabular}
\caption{Bandwidth ceiling vs.\ measured single-stream decode throughput.}
\end{table}

This is the actual explanation for the flattest result anywhere in this project's
benchmark tables: all four concurrency-1 rows, across both engines and both
checkpoints, land between 46.6 and 48.8~tok/s, because every one of them moves the same
$\sim$55~GB over the same memory bus regardless of which server is asking for it. Engine
choice cannot matter until concurrency amortizes those reads across more than one
request at a time, and the remaining 20--25\% of headroom at concurrency~1 can only be
reached by more bandwidth or fewer bytes (quantization), not a different server or a
better flag. This number is approximate by construction --- weight traffic only, no KV
cache reads counted, and the hybrid stack's 48 linear-attention layers do not access
memory the same way the 16 full-attention ones do --- and should be read as a band, not
a precise figure.

\subsection{Serving Traps}

\textbf{vLLM's default \texttt{max\_num\_seqs} will not start this model.}

\begin{lstlisting}
ValueError: max_num_seqs (1024) exceeds available Mamba cache blocks (362).
Each decode sequence requires one Mamba cache block, so CUDA graph capture
cannot proceed.
\end{lstlisting}

This is the hybrid stack asserting itself: the 48 linear-attention layers each need one
Mamba cache block per concurrent sequence, a limit entirely separate from ordinary KV
cache sizing. \texttt{--max-num-seqs 256} fixes it; any copy-pasted serving command for
this model that omits the flag does not work. Separately, the config still carries
\texttt{mrope\_section} and \texttt{mrope\_interleaved} inherited from the multimodal
parent, which vLLM logs as unrecognized rope-parameter keys --- harmless so far, but the
GLM split in Section~\ref{sec:glm} explicitly dropped the same field, and this one
should too.

\subsection{What Removing 1.658\% Actually Buys}

Same H100 80GB, same flags, vLLM~0.27.1:

\begin{table}[h]
\centering
\small
\begin{tabular}{lrrr}
\toprule
 & Parent & Text-only & $\Delta$ \\
\midrule
Model weights & 51.10 GiB & 50.13 GiB & $-$0.97 GiB ($-$1.9\%) \\
KV cache available & 18.27 GiB & 19.84 GiB & \textbf{+1.57 GiB (+8.6\%)} \\
KV cache tokens & 277,435 & 301,465 & \textbf{+8.7\%} \\
Max concurrency @ 32k & 8.47$\times$ & 9.20$\times$ & \textbf{+8.6\%} \\
\bottomrule
\end{tabular}
\caption{What a 1.658\% weight reduction converts into.}
\end{table}

The weight saving itself is small. It converts almost entirely into KV-cache headroom,
which is the number actually worth quoting on a model card rather than the raw
parameter delta.

\subsection{The Benchmark Result That Nearly Got Published Wrong}

512-in / 256-out, greedy decoding, \texttt{ignore\_eos}, one shared OpenAI-compatible
client, H100~80GB.

\begin{table}[h]
\centering
\small
\begin{tabular}{llrrr}
\toprule
Engine & Model & c=1 tok/s & c=16 tok/s & c=16 TTFT p50 \\
\midrule
vLLM 0.27.1 & parent & 47.4 & \textbf{635.8} & 657 ms \\
vLLM 0.27.1 & text-only & 46.6 & 554.7 & 654 ms \\
SGLang 0.5.18 & parent & 48.8 & 225.3 & 10{,}585 ms \\
SGLang 0.5.18 & text-only & 47.2 & \textbf{628.5} & 715 ms \\
\bottomrule
\end{tabular}
\caption{Cross-engine serving throughput, before and after accounting for warmup and planner artifacts.}
\end{table}

Three things to read off this table, in order of how badly the first draft of this
result got them wrong.

\textbf{Warmup moves the number by 21\%.} The first vLLM run on the split measured
457.9~tok/s; the identical server and flags with one discarded warmup pass measured
554.7. Every row above is warmed. An unwarmed row is not comparable to a warmed one, and
mixing the two silently flips conclusions about which checkpoint or engine is faster.

\textbf{``The text-only split is faster'' is false, stated that broadly.} On vLLM the
\emph{parent} is 1.15$\times$ faster at concurrency~16. On SGLang the split is
2.79$\times$ faster --- but only because SGLang's planner allocated the parent
\texttt{max\_running\_requests=5} against the split's 16 and then starved the parent of
concurrency it was never given a chance to use. That is an engine memory-planning
artifact, not a property of either checkpoint.

\textbf{The real result is variance, not speed.} Across the two engines, the parent's
concurrency-16 throughput spans 2.82$\times$ (635.8 vs.\ 225.3); the split's spans only
1.13$\times$ (554.7 vs.\ 628.5) on identical hardware. The split is not faster. It is
\emph{predictable}, and predictability --- not raw throughput --- is what dropping the
vision tower actually buys here. At concurrency~1 all four runs land between 46.6 and
48.8~tok/s regardless of engine or checkpoint: the same 64 text layers do the work in
every case, single-stream decode cannot tell the models apart, and that flatness is
itself the sanity check confirming the surgery did not change the compute path, echoing
Section~\ref{sec:qwen-bandwidth}'s bandwidth-bound explanation for exactly this
pattern.

\section{Case Study: Gemma 4 31B --- Quantization, Not Splitting}
\label{sec:gemma}

Gemma~4 is the model that broke this project's premise. The 12B variant
(\texttt{model\_type: gemma4\_unified}) has no vision \emph{tower} at all --- Google
replaced the 27-layer SigLIP encoder with a single matmul: 48$\times$48 raw pixel
patches through \texttt{LayerNorm} $\to$ \texttt{Linear(6912$\to$3840)} $\to$
\texttt{LayerNorm}, plus factorized coordinate embeddings, and a 40ms-raw-waveform
audio path with no mel spectrogram and no conformer. The whole multimodal path is
52,379,904 parameters: 0.438\% of 11,959,730,224, cross-checked against
\texttt{llama.cpp}'s own \texttt{mmproj-F32} build for this model (209.5~MB $\div$
4~bytes = 52.4M, matching). Removing it saves essentially nothing; quantizing the same
checkpoint to Q4\_K\_M takes 23.9~GB down to 7.1~GB, roughly 160$\times$ the saving
splitting offers. \textbf{Only the 12B is encoder-free} --- E2B, E4B, 26B-A4B, and 31B
are all \texttt{model\_type: gemma4} with conventional SigLIP towers, so checking
\texttt{model\_type} before applying any recipe from this report to a Gemma~4 variant is
not optional.

That realization is what turned this case study from a fourth split into a
quantization project. \texttt{nvidia/Gemma-4-31B-IT-NVFP4} does have a real, removable
27-layer SigLIP tower ($\sim$576M params, $\sim$1.15~GB), but reading the checkpoint's
own quantization metadata shows it is the smallest of three obvious wins, not the
largest:

\begin{table}[h]
\centering
\small
\begin{tabular}{llrr}
\toprule
Component & State & Params & Size \\
\midrule
MLPs (60 layers) & NVFP4, group 16 & 20.8B & 41.6 $\to$ $\sim$11.7 GB \\
Attention (60 layers) & \textbf{untouched, BF16} & 7.93B & 15.85 GB \\
Embeddings (tied = \texttt{lm\_head}) & untouched & 1.41B & 2.82 GB \\
Vision tower (SigLIP, 27 layers) & untouched & $\sim$576M & $\sim$1.15 GB \\
\bottomrule
\end{tabular}
\caption{NVIDIA's \texttt{Gemma-4-31B-IT-NVFP4}, by component ($\sim$31.5 GB against a
real 32.63 GB, consistent with the accounting).}
\end{table}

NVFP4 advertises 4-bit weights \emph{and} activations, but the checkpoint's
\texttt{hf\_quant\_config.json} excludes 63 modules, 60 of them the attention blocks,
listed one layer at a time. Only the MLPs were quantized, leaving unquantized attention
as the single largest block in the file. Leaving it BF16 is a defensible choice ---
NVFP4 quantizes activations too, and attention activations are exactly the
outlier-heavy case SmoothQuant and AWQ exist to handle --- but it is also a deviation
from NVIDIA's own practice elsewhere: their DeepSeek-V4 and Qwen3-Omni builds ship FP4
experts alongside FP8 attention, and \texttt{mtq.NVFP4\_MLP\_WEIGHT\_ONLY\_CFG} is a
named recipe in ModelOpt, so MLP-only is a deliberate configuration rather than an
oversight. This project followed three changes from that observation, in order: strip
the vision tower (pure deletion, exactly verifiable), take attention weights to FP8 (the
real size win), and pair the unpaired MTP drafter (free throughput, and an open
question). The outcome: \textbf{32.63~GB~$\to$~23.00~GB (--29.5\%)}, NVIDIA's NVFP4 MLPs
copied verbatim, 230 attention projections requantized to FP8, no vision tower,
retargeted onto \texttt{Gemma4ForCausalLM}, with generation quality preserved. This
section is deliberately written as a narrative rather than a result sheet, because most
of the value in building it was in the wrong turns.

\subsection{Decision: Weights-Only FP8, Not W8A8}

FP8 has two independent halves with very different risk profiles. The weight half is a
pure function of the stored weights --- no calibration set, no forward passes, nothing to
get wrong beyond scale arithmetic. The activation half needs calibration data and is
exactly where attention's outliers cause trouble. Weights-only still takes attention
from 15.85~GB to 7.93~GB and still speeds up batch-1 decode, since decode at batch~1 is
memory-bandwidth-bound rather than compute-bound; what it gives up is a prefill win,
because BF16 activations cannot feed FP8 tensor cores. This project chose weights-only
because it is the version that definitely ships, and because the headline measurement in
Section~\ref{sec:gemma-headline} wanted exactly one variable moving at a time. A
calibrated W8A8 variant remains a sensible follow-up, not a rejected idea.

\subsection{Decision: Modify NVIDIA's Checkpoint, Not Requantize from BF16}

The alternative was starting from \texttt{google/gemma-4-31B-it} and running one
mixed-precision ModelOpt pass from scratch. Rejected: it needs a 62~GB download and a
calibration pass, and the result would no longer be comparable to NVIDIA's own published
benchmark numbers, since any accuracy delta could then come from this project's
calibration set rather than from the FP8 attention change under test. Working in place
on NVIDIA's own checkpoint keeps their calibrated NVFP4 MLPs bit-identical, so exactly
one thing changes between the baseline and the result.

\subsection{What Did Not Work: Loading the Checkpoint at All}

\texttt{transformers} cannot load this checkpoint. \texttt{quant\_method: modelopt} is
not among the methods it supports (\texttt{fp8}, \texttt{mxfp4},
\texttt{compressed-tensors}, and \texttt{quark} are; \texttt{modelopt} is not), and it
does not error on the unknown method --- it ignores the quantization config entirely,
loads the packed \texttt{uint8} tensors as if they were dense weights, and dies on a
shape mismatch:

\begin{lstlisting}
ckpt: torch.Size([21504, 2688])  vs  model: torch.Size([21504, 5376])
\end{lstlisting}

Exactly 2$\times$ on the input dimension, which is itself the tell that NVFP4 packs two
4-bit values per byte, and the clue that made the unpacker straightforward to write once
noticed.

\subsection{The Fix That Unblocked Everything: Verify by Expansion}

Weight-only quantization is exactly simulatable. An FP4 or FP8 kernel does not multiply
in 4 or 8 bits; it dequantizes the stored weight to a wide type first and multiplies
against the activations in that wide type. So a checkpoint whose weights have been
round-tripped through the format computes bit-identical results to the real deployed
kernel, and accuracy can be measured on an expanded BF16 copy without ever running the
quantized kernel itself. Two consequences followed directly: \textbf{NVFP4 accuracy work
does not need Blackwell hardware}, only throughput measurement does, which overturned
this project's earlier plan (that plan had assumed scoring required SM100); and the
missing \texttt{transformers} loader became a non-issue, since expansion sidesteps it
entirely. What expansion does \emph{not} simulate is NVFP4's activation quantization, so
an expanded NVFP4 checkpoint is very slightly optimistic relative to a true end-to-end
deployment. Both sides of every comparison in this section are expanded the same way, so
the comparison stays fair for the specific weight change under test, but none of the
numbers below are a claim about deployed accuracy on real NVFP4 tensor cores.

NVFP4's on-disk layout, confirmed against the real tensors rather than assumed from the
spec:

\begin{lstlisting}
weight          [out, in/2]   uint8       two e2m1 values per byte, low nibble first
weight_scale    [out, in/16]  float8_e4m3 one per group of 16
weight_scale_2  scalar        float32     one per tensor
w = e2m1(nibble) * weight_scale * weight_scale_2
\end{lstlisting}

Two checks run before trusting any downstream number from the dequantizer:
\texttt{model.layers.\{0,30,59\}.self\_attn.q\_proj.weight} must be bit-identical
(\texttt{max|d| = 0.0000}), since attention was never quantized by NVIDIA and must
survive stripping and renaming untouched; and \texttt{weight\_scale.max()} $\times$
\texttt{weight\_scale\_2} $\times$ 6 must reproduce the original tensor's absmax exactly
(measured: 0.2617), the cheapest possible identity check that the largest e2m1
magnitude times the largest block scale reproduces the original range. Run both before
trusting anything else this section reports.

\subsection{Five Bugs the Real Checkpoint Found That Local Tests Did Not}

\begin{enumerate}[leftmargin=1.5em]
\item \textbf{\texttt{*self\_attn*} matches \texttt{q\_norm} and \texttt{k\_norm}.} An
  early run quantized all 120 of them. They are 1-D RMSNorm scales --- no size to win by
  quantizing them, and among the most precision-sensitive weights in the model. Caught
  only because the converted tensor count came to 350 where projections should total
  230; fixed by targeting \texttt{*\_proj} specifically and structurally refusing any
  1-D tensor, since the pattern had already proven easy to get wrong once.
\item \textbf{The quantization schema is a hybrid.} \texttt{quantization\_config}
  carries modelopt keys (\texttt{quant\_algo}, \texttt{quant\_method}) and
  compressed-tensors keys (\texttt{config\_groups}, \texttt{ignore},
  \texttt{kv\_cache\_scheme}) in one block. Reading only the modelopt half silently
  loses the group size, which is nested inside \texttt{config\_groups}, and the KV
  cache format, which is named differently in each half.
\item \textbf{Tied embeddings defeat name matching.} The exclusion list names
  \texttt{lm\_head}, but the tensor ships under the key \texttt{embed\_tokens}, so
  matching on names alone reports BF16 weights as already quantized. The dtype recorded
  in the tensor header is the authority; the config is only a claim about it.
\item \textbf{Renaming a prefix invalidates the exclusion patterns.} Moving
  \texttt{model.language\_model.*} to \texttt{model.*} leaves every \texttt{ignore}
  entry pointing at a path that no longer exists after the rename, so they match
  nothing, and the next reader of the checkpoint concludes attention was quantized when
  it was not.
\item \textbf{Patterns matched full tensor names, not modules.} Any user-supplied
  pattern without an explicit \texttt{.weight} suffix would silently convert nothing at
  all, passing without any visible error.
\end{enumerate}

\texttt{230}, incidentally, is not a typo for 240: Gemma~4 sets
\texttt{attention\_k\_eq\_v}, so ten layers share K and V and have no separate
\texttt{v\_proj} --- worth checking the architecture before assuming a tensor count is
simply wrong.

\subsection{What Did Not Work: Scoring on Raw Prompts}
\label{sec:gemma-rawprompts}

The first scoring run produced clean-looking numbers claiming this project's checkpoint
\emph{beat} NVIDIA's own build, and it passed the ``no worse than baseline'' gate
outright. It was entirely noise. The tell was in the generations themselves: the BF16
\emph{reference} model was emitting \texttt{France is France is France is\ldots} The
control experiment that cleared it: the untouched original produces the identical
degenerate output on that exact prompt, and through the chat template both the original
and this project's copy correctly answer \texttt{Paris}. So the bug was never in the
extraction --- \texttt{gemma-4-31B-it} is instruction-tuned, and the scoring script was
feeding it raw completion prompts rather than the chat template it was trained to expect.
Scoring in that degenerate regime inflates divergence by roughly an order of magnitude
in arbitrary directions, which is precisely how a \emph{more} quantized checkpoint
managed to come out looking closer to the reference than a less-quantized one did.
\texttt{agreement} grew a \texttt{--chat} mode and a fixed \texttt{CHAT\_PROMPTS} set as
a direct result; \texttt{verify} kept its raw prompts unchanged, because an extraction
must stay bit-exact on any input, degenerate or not, which is a different claim from
scoring quantization damage and belongs to the other bar described in
Section~\ref{sec:methodology}. The general lesson generalizes past this one incident: a
plausible result that favors the work already done is the one to distrust most.

\subsection{Do Not Read KL as Benchmark Accuracy}

A KL divergence of 0.145 looks alarming set next to NVIDIA's own published GPQA numbers,
85.35\% vs.\ 85.80\%, which \emph{is} genuinely near-lossless. There is no contradiction
between the two: KL computed over a 262k-token vocabulary is far more sensitive than a
multiple-choice accuracy metric, and a 9.5\% relative weight error at 4 bits can leave
benchmark scores essentially unchanged while still producing a KL number that looks
large in isolation. This nearly triggered a false alarm about a bug in the dequantizer
before the comparison was reframed correctly: judge one quantization against another
quantization, never against an absolute divergence number with no reference point.

Measured through the chat template against the BF16 reference:

\begin{table}[h]
\centering
\small
\begin{tabular}{lrr}
\toprule
 & Worst top-1 & Worst KL \\
\midrule
B --- NVIDIA NVFP4 & 0.7826 & 1.448e-01 \\
C --- ours, +FP8 attention & 0.9130 & 1.716e-01 \\
\bottomrule
\end{tabular}
\caption{Divergence from the BF16 reference, NVIDIA's build vs.\ this project's build.}
\end{table}

C adds 1.19$\times$ baseline's KL divergence. Its top-1 agreement came out
\emph{higher} than B's on all six test prompts, which should not happen for a strictly
more quantized model --- six prompts and roughly 150 token positions is a small enough
sample that this pattern is plausibly the sample size talking, not a real effect. Both
rows should be read as indicative, not final, until re-run on a larger prompt set.

\subsection{The Headline: Quantizing the Target Costs Speculative Acceptance}
\label{sec:gemma-headline}

The open question this investigation was actually built to answer. The drafter model
(\texttt{google/gemma-4-31B-it-assistant}, 0.47B parameters) was trained against the
BF16 target; quantizing the target moves its output distribution away from what the
drafter learned to predict. Speculative decoding stays exact by construction regardless
--- the target verifies every drafted token and discards whatever it would not have
produced itself --- so this cannot change output quality. It changes throughput, and
under exact speculative decoding throughput is entirely a function of acceptance rate.
Measured as tokens accepted per target forward pass --- the arithmetic ceiling on
speedup, and a property of the two distributions rather than of any kernel implementing
them:

\begin{table}[h]
\centering
\small
\begin{tabular}{lrr}
\toprule
Target & tok/fwd & vs.\ BF16 \\
\midrule
R --- base BF16 & 2.220 & --- \\
B --- NVIDIA NVFP4 & 1.655 & $-$25.4\% \\
C --- ours, +FP8 attention & 1.247 & $-$43.8\% \\
\bottomrule
\end{tabular}
\caption{Speculative-decoding acceptance, tokens per target forward pass, by how much
the target is quantized.}
\end{table}

\textbf{Monotonic in the amount of quantization applied}, which is what makes this
result credible where the top-1 numbers in Section~\ref{sec:gemma-rawprompts} were not.
A published ``up to 3$\times$'' speedup claim for the BF16 pairing does not transfer to
a quantized target: NVIDIA's own build already surrenders a quarter of the benefit
before this project's FP8 attention change adds anything on top of it.

Three caveats that matter to how far this generalizes. The wall-clock column should be
ignored entirely: the BF16 reference measured a 0.97$\times$ ``speedup'' ---
speculative decoding running \emph{slower} than plain generation --- because
\texttt{transformers}'s assisted-generation overhead dominates at this model scale, a
property of the measurement harness, not of either model, and exactly why tokens per
forward pass is the number reported instead. The measurement used
\texttt{assistant\_model} draft-model speculation, not necessarily the native MTP path
the drafter ships for. And it covers four prompts at 96 new tokens each --- directionally
clear, not a precise coefficient to quote past one decimal place. The practical reading
stands regardless of those caveats: if throughput is the reason to want speculative
decoding at all, quantizing the target is not free, and the trade against 30\% smaller
weights is a real cost to weigh, not an assumption to make.

\subsection{The KV Cache, and a 2$\times$ Error in the First Estimate}

The initial analysis assumed uniform layers and claimed 80.0~KB/token, 21.47~GB at full
context --- ``as large as the weights.'' That was wrong, and it was wrong because
Gemma~4's layers are heterogeneous, a fact that surfaced only because
\texttt{transformers} raises \texttt{AmbiguousGlobalPerLayerAttributeError} on a global
read of \texttt{num\_key\_value\_heads} rather than silently returning a number for the
wrong layer type:

\begin{lstlisting}
10 x full_attention     kv_heads=4   head_dim=512   -> 40.0 KB/token
50 x sliding_attention  kv_heads=16  head_dim=256   -> capped at ~0.4 GB by the window
\end{lstlisting}

The counter-intuitive part is that the sliding layers carry four times the KV heads of
the global ones, but their 1024-token window caps them at a fixed $\sim$0.4~GB
regardless of how long the context runs --- only the ten global layers actually grow with
context length: 1.75~GB at 32,768 tokens, 5.79~GB at 131,072, 11.16~GB at the full
262,144-token window. So at the full window the cache is roughly half the quantized
weights, not equal to them as the first, uniform-layer estimate claimed. Worth knowing
before optimizing further, but it does not outrank the weight savings reported above,
and the stronger original claim was purely an artifact of assuming a homogeneous model
where none exists.

\subsection{Two Things Reached the Published Repo Before Being Caught}

Both are the same bug wearing different clothes: metadata that describes the model
\emph{before} the surgery, carried forward as if it still applied after. \textbf{Loose
files travel with a checkpoint.} The copy step that carried non-safetensors files
forward brought NVIDIA's own \texttt{README.md} (describing their build, not this
project's result) and a 2,155-line ModelOpt quantizer dump listing the vision tower and
reporting attention as unquantized --- both landed in the upload, and neither was true of
the actual result shipped. There is now a \texttt{STALE\_METADATA} skip set, with a test
asserting the tokenizer and chat template still travel with the checkpoint, so the guard
does not become an excuse to drop files the model genuinely needs to load.
\textbf{Promoting a nested config exposes fields to new validators.} Gemma~4's
\texttt{text\_config} carries \texttt{num\_experts: null},
\texttt{expert\_intermediate\_size: null}, and \texttt{top\_k\_experts: null}, all inert
because \texttt{enable\_moe\_block} is \texttt{False}. Nested inside \texttt{text\_config}
nobody checks them; promoting \texttt{text\_config} to the top level exposes them to the
Hugging Face config parser, which rejects \texttt{"num\_experts" must be a number}. The
fix --- dropping the three inert keys --- was only trusted once \texttt{Gemma4TextConfig}
was constructed both ways and every resolved field was asserted to match; dropping
config keys because they ``look unused'' is otherwise a guess, not a verified fix.

\subsection{vLLM and This Checkpoint: Three Wrong Answers Before the Right One}
\label{sec:gemma-vllm}

The published model card told people to run \texttt{vllm serve --quantization modelopt}.
That instruction was written from what the format's documentation implies, never
actually run, and it does not work. Correcting it took two rented GPU instances and
about \$7, mostly because the first two candidate explanations were both wrong and both
\emph{felt} conclusive at the time. The final answer: this checkpoint cannot be served
by vLLM at all today, for a reason that has nothing to do with Gemma~4 or with vLLM's
Gemma support specifically, and everything to do with the weights-only FP8 decision
recorded earlier in this section.

\subsubsection*{The first wrong answer: ``vLLM doesn't support Gemma 4''}
The observed error was \texttt{AmbiguousGlobalPerLayerAttributeError: 'head\_dim' is a
per-layer attribute}, and the conclusion drawn from it was an upstream architecture
support gap. That conclusion cost nothing to check and was false: a file listing of the
release tag --- no GPU, no download, one \texttt{curl} call --- shows \texttt{gemma4.py},
\texttt{gemma4\_mm.py}, and \texttt{gemma4\_mtp.py} all present in v0.27.1, with
\texttt{git log} dating the first of them four months before that release. \emph{Before
renting a GPU to explain a failure, check whether the cheapest possible observation
already falsifies the explanation.}

\subsubsection*{The second wrong answer: ``it's a transformers version skew''}
Falsifying a cause is not the same as finding one. The next hypothesis came from two
pieces of evidence inside the error itself: the exception's home library is
\texttt{transformers}, not vLLM, which is evidence about which side of the stack
changed; and vLLM~0.27.1 pins \texttt{transformers>=5.5.3} with no upper bound while
unfuse requires \texttt{transformers>=5.15.1}, so vLLM was necessarily running against a
transformers version newer than anything it had been tested with. That hypothesis is
testable on the config alone, with no weights and no GPU needed, and a version sweep
bisected it exactly:

\begin{lstlisting}
5.5.3, 5.8.0, 5.11.0, 5.13.1, 5.14.0, 5.14.1   head_dim = 256    global read OK
5.15.0, 5.15.1                                 AmbiguousGlobalPerLayerAttributeError
\end{lstlisting}

Pinning \texttt{transformers==5.14.1} does clear the config error, and that is a real,
necessary finding affecting any Gemma~4 31B on vLLM, Google's own build included --- but
it was not the actual blocker. It only cleared the way to the next failure, an
\texttt{AssertionError} deep inside the QKV loader. \emph{Fixing the first error tells
you what the second error is, not that you were right about the cause.} The
\texttt{transformers} guard turned out to be correct, not a regression: Gemma~4 31B
genuinely has two head dimensions, 256 on the 50 sliding layers and 512 on the 10
full-attention layers, so a global \texttt{head\_dim} read is meaningless by
construction, and vLLM~0.27.1 reading it globally returns 512 --- wrong for 50 of 60
layers.

\subsubsection*{The control experiment that inverted the conclusion}
The question worth answering was not ``why does this fail'' but ``is this ours?'' The
control is running the unmodified upstream artifact through the identical serving path:

\begin{lstlisting}
nvidia/Gemma-4-31B-IT-NVFP4                       ->  SERVED OK
sartajbhuvaji/Gemma-4-31B-IT-NVFP4-FP8attn-text   ->  AssertionError in load_qkv_weight
\end{lstlisting}

Upstream serves, on an A100, not even on Blackwell hardware. So vLLM can serve Gemma~4
31B, the transformers pin is sufficient for that, and the blocker is something this
project's build introduced, not an upstream gap. \emph{A control run on the unmodified
upstream artifact is what separates ``their bug'' from ``ours,'' and it is worth its
cost early rather than late.}

\subsubsection*{Instrumenting instead of guessing}
The assertion carried no shapes in its message, so the loader was monkeypatched to
print them at the exact point of failure rather than reasoning about what they might be
from the stack trace alone:

\begin{lstlisting}
shard=k   param=(16384, 2688)   loaded=(4096, 5376)
\end{lstlisting}

$2688 = 5376/2$ is the tell --- the same 2$\times$ factor that identifies FP4 packing
everywhere else in this section. vLLM had allocated the attention parameter as if it
were NVFP4-packed, while the stored tensor is actually full-width FP8: vLLM was applying
NVFP4 to attention, meaning it had never read the \texttt{quantized\_layers} mapping
that says otherwise. Reading the dispatch source rather than the documentation found
why --- two exact-match gates, neither visible from the config file alone:
\texttt{quant\_algo} is literally the string \texttt{"NVFP4"}, and vLLM's mixed-precision
path is reachable only when that string equals \texttt{"MIXED\_PRECISION"} exactly, so it
was never entered; and \texttt{quantized\_layers} keys in the checkpoint are globs like
\texttt{model.layers.0.self\_attn*}, but the lookup vLLM performs is exact dict
membership with no glob expansion, so they match nothing. The correct form names each
projection individually: 230 FP8 entries beside 180 NVFP4 MLP entries. \emph{When a
config key appears to be silently ignored, read the dispatch code --- the gate is often a
literal string comparison, and no amount of plausible-looking JSON will pass it.}

\subsubsection*{The real blocker, and why it is the worst kind}
With both defects above patched, the model loaded, ran, and generated fluent nonsense:

\begin{lstlisting}
>>> What is the capital of France? Answer in one word.
same dePK singleale dePK singleale dePK singleale ...
\end{lstlisting}

\texttt{ModelOptFp8LinearMethod.create\_weights} registers \texttt{input\_scale} as a
required per-tensor parameter, and \texttt{process\_weights\_after\_loading} calls
\texttt{layer.input\_scale.max()}. vLLM's ModelOpt FP8 path is W8A8 only. This checkpoint
is deliberately W8A16 --- FP8 weights, BF16 activations, chosen for exactly the accuracy
reasons recorded earlier --- so there is no \texttt{input\_scale} tensor to load, and the
scale stays at its initialization sentinel. Nothing errors, because FP8 and BF16 weights
have identical tensor shapes: every shape check in the loader passes, and only the
arithmetic downstream of it is wrong. This is strictly worse than the crash it replaced
--- a crash self-reports, and a checkpoint that serves confident garbage does not.
\emph{Prefer the failure mode that announces itself: a dtype swap that preserves tensor
shape will pass every structural check a loader makes, and fail silently past all of
them.}

So the decision that is correct on accuracy grounds (Section~\ref{sec:methodology}'s
earlier discussion) is exactly the one with no serving path today. Fixing the two
dispatch defects in isolation would convert a loud failure into a silent one, which is
worse, not better --- they are worth fixing for metadata correctness regardless, but
\texttt{quant\_algo} should stay \texttt{NVFP4} until either a calibrated W8A8 variant
exists or vLLM gains FP8 weight-only support in the ModelOpt path specifically.

\subsubsection*{What actually runs, tested}
Stock released vLLM, not a nightly build --- ``install the nightly'' is a materially
different instruction to publish on a model card than a pinned release:

\begin{lstlisting}
pip install vllm==0.27.1 'transformers==5.14.1'
python tools/dequantize.py ./ckpt ./ckpt-bf16
vllm serve ./ckpt-bf16
\end{lstlisting}

On an H100 PCIe 80~GB: engine up in 103~s, 36.5~tok/s on a batch of 5, eager mode.
Correct answers on every test prompt: \texttt{Paris}; \texttt{408} with correct working
shown; \texttt{Das Wetter ist heute kalt.}; Rayleigh scattering explained correctly;
\texttt{my\_string[::-1]}. A vLLM nightly serves the same expansion at 42~s and
48.8~tok/s. Expansion costs back the size saving --- 23~GB becomes 61~GB --- so it needs
an 80~GB card regardless of which vLLM build serves it.

\subsubsection*{Corrections to claims this project had already published}
Three claims in earlier versions of this project's own documentation turned out to be
wrong, corrected only once the vLLM investigation above forced a closer look.
\textbf{``NVFP4 executes only on Blackwell (SM100+)'' is false under vLLM}: it selects
\texttt{MarlinNvFp4LinearKernel}, and \texttt{ModelOptMixedPrecisionConfig.get\_\allowbreak
min\_capability()} reports~75, with a comment in the source recording validation on a T4
(SM75) and an A100 (SM80). Blackwell provides native FP4 tensor cores; without them,
Marlin repacks to W4A16 at load time, so the \emph{stored} size saving is not
automatically a \emph{runtime} memory saving --- a 40~GB A100 ran out of memory loading
this 23~GB checkpoint for exactly this reason. \textbf{The FP8 KV cache needs SM89+, not
Blackwell}; below that, vLLM refuses to start with \texttt{FP8 KV cache is not supported
by the Triton attention backend} and names its own fix, \texttt{--kv-cache-dtype
bfloat16}. And \textbf{vLLM~0.27.1 has no \texttt{quantized\_layers} support at all} ---
it parses only \texttt{exclude\_modules} (the legacy \texttt{hf\_quant\_config.json}
field) or \texttt{ignore}. NVIDIA's own build serves correctly on this release
specifically because its 60 attention blocks sit in \texttt{exclude\_modules}, which
0.27.1 does read; mixed precision as this project configured it cannot be expressed to
that vLLM release in any form.

\subsubsection*{The meta-lesson}
Two of the three explanations along the way were wrong, and each one was written down
with real confidence before the next observation demolished it, including in this
project's own earlier published documentation. The pattern that eventually worked: the
cheapest falsifying observation first, then a control run against the unmodified
upstream artifact, then instrumentation at the exact point of failure, then the serving
library's own dispatch source --- and only after all four, a claim, tested against the
exact software versions named in it.

\section{Cross-Cutting Lessons}
\label{sec:lessons}

The three case studies above each have their own narrative; this section pulls out the
lessons that generalize past any one of them, grouped by where they actually bite rather
than by which model surfaced them first. Each cost real time or real GPU-hours to find.

\subsection{Config and Metadata Traps}

\textbf{A config field can describe the model before the surgery, not the model that
actually ships.} This is the single most common failure pattern across every case study
in this report: \texttt{deepstack\_visual\_indexes} in Section~\ref{sec:qwen} looked
like a problem and was not; the exclusion-list name mismatch and the invalidated
\texttt{ignore} prefixes in Section~\ref{sec:gemma} both looked like correct
configuration and were not; the loose \texttt{README.md} and quantizer dump that
traveled into a published repo in Section~\ref{sec:gemma} described NVIDIA's build, not
this project's result. The fix is the same every time: verify against the tensor
headers and the dispatch code, not against what the config claims.

\textbf{A class that does not populate \texttt{architectures} leaves nothing for
\texttt{AutoModelForCausalLM} to dispatch to.} True of \texttt{Glm4Config}
(Section~\ref{sec:glm}) and of \texttt{Qwen3\_5TextConfig} (Section~\ref{sec:qwen}),
independently, on two unrelated model families. Setting \texttt{architectures}
explicitly after reconstructing a config is not an edge case to remember occasionally;
treat it as a required step of any split.

\textbf{\texttt{total\_size} in \texttt{model.safetensors.index.json} is not always
bytes.} For GLM-4.6V-Flash it reports the parameter count instead (Section~\ref{sec:glm}).
Verify against real file sizes rather than trusting the field name.

\subsection{Serving-Engine Quirks}

\textbf{Each engine's own benchmark tool disagrees with every other one's.}
\texttt{vllm bench serve} and \texttt{sglang.bench\_serving} sample prompts differently,
warm up differently, and define ``duration'' differently. Run the Qwen3.8-27B comparison
in Section~\ref{sec:qwen} against each tool's own harness and SGLang appeared to beat
vLLM on time-to-first-token by 3$\times$; run both against one shared OpenAI-compatible
client over plain HTTP instead, and vLLM won by 1.7$\times$. Cross-engine numbers are
only meaningful from a single shared harness. Pin output length with \texttt{ignore\_eos}
as well, or a chattier model's verbosity becomes the throughput measurement instead of
the thing actually under test.

\textbf{NVFP4 does not require Blackwell, and FP8 KV cache does not require it
either, at a different threshold.} vLLM selects \texttt{MarlinNvFp4LinearKernel} and
runs it from SM75 up; without Blackwell's native FP4 tensor cores, Marlin repacks to
W4A16 at load time, so a checkpoint's \emph{stored} size saving is not automatically a
\emph{runtime} memory saving, and a 23~GB NVFP4 checkpoint OOM'd a 40~GB A100 for
exactly this reason (Section~\ref{sec:gemma-vllm}). The FP8 KV cache separately needs
SM89+; below that, pass \texttt{--kv-cache-dtype bfloat16}.

\textbf{Installing a tool into a venv does not put it on \texttt{PATH}.} vLLM shells
out to \texttt{ninja} at build time and fails with a bare \texttt{FileNotFoundError}
surfaced three stack frames deep as \texttt{Engine core initialization failed}.
Activating the venv puts \texttt{\$VENV/bin} on \texttt{PATH}; running
\texttt{\$VENV/bin/vllm} directly does not, so a correctly installed \texttt{ninja}
stays invisible to the engine that needs it. Separately, vLLM and SGLang pin different
torch builds, and installing both into one venv lets the second silently downgrade the
first --- always give them separate environments.

\textbf{The reasoning field has two names.} vLLM~0.27.1 returns a model's thinking
trace as \texttt{message.reasoning}; SGLang returns the same concept as
\texttt{message.reasoning\_content}. Client code written against one silently drops it
on the other with no error.

\textbf{A daemon thread holding CUDA can abort the process at exit, after every check
has already passed.} Streaming generation handed to a \texttt{Thread(...,
daemon=True)} can still be inside \texttt{generate()} when the interpreter tears down
CUDA at exit, producing \texttt{terminate called without an active exception} and a
non-zero exit code after every metric has already been computed and written to disk.
Non-daemon, plus an explicit \texttt{join()}, avoids it.

\subsection{Infrastructure and Process Traps}

\textbf{A background job's own echo is not evidence it started.} \texttt{pip install
huggingface\_hub[cli]} warns that it installed to \texttt{\textasciitilde/.local/bin},
which is not on \texttt{PATH} on a fresh cloud image. A backgrounded
\texttt{nohup hf download ... \&} then fails with \texttt{nohup: failed to run command
'hf'} \emph{inside the log file}, while the launching shell happily reports the download
as started. Twenty minutes of a rented GPU instance went to a download that never began.
Use absolute paths for anything backgrounded, and verify bytes actually landing on disk
rather than trusting a shell's own echo that a command launched.

\textbf{\texttt{pkill -f <pattern>} can match your own SSH session.} The remote
command string lives in the remote shell's own \texttt{argv}, so \texttt{pkill -f
pipeline.sh} issued over SSH can kill the shell that is running the SSH session itself
--- connection drops, exit 255, and whatever was mid-flight dies with it. This happened
three times before being tracked down. The \texttt{[p]attern} bracket trick only helps
when the pattern string is not also present verbatim somewhere else in the full command
line; a heredoc containing the same string defeats it too. Kill by PID instead, or write
the command to a file and \texttt{scp} it over rather than passing it inline.

\textbf{\texttt{hf upload} creates the destination repo public by default, with no
prompt.} A 54.64~GB checkpoint began publishing to a public repository before its model
card existed to explain what it was. Create the destination repo explicitly ahead of
time, or check \texttt{model\_info(...).private} immediately after upload starts and
flip it if it is wrong.

\textbf{A backgrounded job that must outlive the SSH session needs
\texttt{setsid}.} Plain \texttt{nohup ... \&} still dies when its parent session is
killed; \texttt{setsid nohup ... < /dev/null \& disown} survives a dropped connection.

\textbf{A \texttt{.env} file edited on Windows silently breaks every API key read from
it.} Files under \texttt{/mnt/c} in WSL carry CRLF line endings; sourcing the file
leaves a trailing \texttt{\textbackslash r} on the last value read, and the service
reports \texttt{invalid-api-key} for a key that is, character for character, correct.
\texttt{file .env} reports \texttt{CRLF line terminators} directly;
\texttt{sed -i 's/\textbackslash r\$//' .env} fixes it, and the launch tooling now
strips CRLF on every read rather than trusting the file to already be clean.

\textbf{\texttt{--system-site-packages} mixes apt and pip packages badly.} A venv
built with that flag inherited apt's torchvision~0.22 while pip separately installed
torch~2.13, so \texttt{torchvision::nms} failed to register --- surfacing as a
completely unrelated-looking \texttt{ModuleNotFoundError: Could not import module
'Gemma4Config'}. pip then skipped installing its own torchvision entirely, since apt's
copy already ``satisfied'' the version requirement, and apt's scipy then collided with
the venv's numpy on top of that. The fix was a clean venv with no system packages and
torch plus torchvision installed from the same CUDA wheel index together.

\textbf{A shell script without a shebang runs under \texttt{/bin/sh}, not bash.} A
setup script used \texttt{exec > >(tee ...)}, a bash-only process substitution; under
\texttt{sh} it dies instantly and leaves no log at all, because the redirect that would
have recorded the error was itself the thing that failed.

\textbf{Lambda's \texttt{POST /ssh-keys} with a null \texttt{public\_key} does not
error.} It silently generates its own keypair instead and returns the private half in
the API response body, which nobody necessarily saves. \texttt{ssh\_key\_names} is fixed
at instance launch time, so discovering this after launch means terminate and relaunch,
not a fix in place. The root cause was a shell bug, not an API quirk:
\texttt{python3 -c "..." P="\$V"} passes \texttt{P=} as a script \emph{argument}, not an
environment variable, so the key content silently never reached the script that was
supposed to read it. Environment assignments must precede the command they apply to.
This project's launch script now verifies the registered key against the local
\texttt{.pub} file before ever launching an instance against it, rather than trusting
that registration worked.

\section{Conclusion and Future Work}
\label{sec:conclusion}

Three models, three different reasons to split or not split them. GLM-4.6V-Flash
(Section~\ref{sec:glm}) is the clean case: a genuinely separable vision tower, a
bit-exact split, and a memory-win claim that needed correcting once the parent was
actually benchmarked rather than only the derivative. Qwen3.8-27B
(Section~\ref{sec:qwen}) is popular enough that packaging, not size, is the whole
argument for splitting it at all --- its tower is 1.658\% of the checkpoint. Gemma~4 31B
(Section~\ref{sec:gemma}) is the model that rejected the premise outright: its
``tower'' is a single matmul, splitting saves almost nothing, and the actual
investigation became quantization, with the speculative-decoding acceptance result in
Section~\ref{sec:gemma-headline} as the one measurement in this report that was not
really about splitting at all. The throughline across all three is not a method; it is
a discipline --- \texttt{unfuse verify}'s bit-exact bar for operations that should not
change arithmetic, \texttt{unfuse agreement}'s comparative bar for operations that
necessarily do, and a running habit of distrusting any result that happens to favor the
work already done, demonstrated concretely in Sections~\ref{sec:methodology} and
\ref{sec:gemma-rawprompts}, independently, on two different models.

\subsection{Why Not to Prune for Size Instead}

Worth stating directly, since a reasonable alternative to splitting is pruning the
language model itself. The literature argues against it here. At 15\% backbone
compression, structural pruning of vision-language models degrades multimodal
performance by 9.13\% while text-only performance drops only 1.07\%
\citep{lvlmpruning2026}: vision-language alignment is the first thing to degrade, and it
degrades without being specifically targeted, so hunting for vision-specific neurons to
prune surgically is the wrong level of precision --- width-pruning generically and
letting vision-serving capacity die along with it is both simpler and consistent with
what actually happens. \citet{minitron2024} remains the credible path to a genuinely
smaller model instead of a repackaged one, having demonstrated exactly a 12B$\to$8B
reduction (Mistral NeMo to MN-Minitron-8B, MMLU 69.5 against the teacher's own 68.0),
but width-pruning only, never depth --- layer-dropping destroys GSM8K and HumanEval while
leaving MMLU looking deceptively unaffected. A $\leq$30\% target is the credible range,
with 45\% the measured tipping point, at a real cost of roughly 18k H100-hours at the
94B-token retraining tier ($\sim$\$40k), or \$150--300k for a full from-scratch
replication --- not a project-weekend undertaking. And multimodal adaptation specifically
suppresses text ability in ways that do not reverse cheaply:
\citet{linguadistill2026} documents a loss on text-only tasks that is ``hard to undo
with further fine-tuning,'' with existing remedies needing added adapters or alignment
modules rather than a lightweight pass. There is no dormant text-only model sitting
underneath a multimodal one's weights, waiting to be fine-tuned back out. Splitting a
model that already has a separable encoder gets the model this project actually wants
without any of that cost; pruning one that does not is a different, much more expensive
project.

\subsection{What Is Still Open}

\textbf{The MINER ablation nobody has run on an encoder-free model.}
\citet{miner2024} found that roughly 2\% of FFN neurons become modality-specific in
encoder-\emph{based} multimodal models, concentrated in shallow layers, and that
deactivating that 2\% collapses multimodal performance (Qwen2-VL: 0.56~$\to$~0.24). It
never ran the converse direction --- ablate the vision-specific neurons, then measure
\emph{text-only} performance, to see whether any serving capacity for vision was
quietly living in the shared FFN all along --- and nobody has run either direction on an
encoder-free model. Every published modality-pruning result in this space studies
LLaVA, InternVL, or Bunny, all of which have separable vision towers; Gemma~4 12B
(Section~\ref{sec:gemma}) does not, by construction. The script for this exists, scores
FFN neurons by activation ratio on image vs.\ text tokens, zeroes the top slice, and
watches both NLLs move. If vision NLL rises while text NLL holds steady, vision-serving
capacity emerged in the shared FFN despite the encoder-free architecture, which would
mean this report's 0.438\% ceiling for Gemma~4 12B's removable multimodal footprint is
the wrong number to quote. It has never been run, and is kept outside this project's
repository until it produces a result worth reporting, following the same discipline
applied to every other number in this report. Estimated cost: roughly one GPU-day.

\textbf{The FP8-attention serving path.} Section~\ref{sec:gemma-vllm} leaves this
project's best Gemma~4 31B checkpoint correct on accuracy grounds and unservable by any
released vLLM build, because the W8A16 choice that is right for accuracy has no
supporting kernel path in vLLM's current ModelOpt integration. Two ways forward exist,
and neither has been attempted yet: calibrate a true W8A8 variant, accepting the
activation-quantization risk this report deliberately avoided, or wait for upstream FP8
weight-only support in vLLM's ModelOpt path. Until one of them lands, the published
checkpoint runs only through \texttt{tools/dequantize.py} back to BF16, at the memory
cost that implies.

\textbf{The MTP pairing, verified only from metadata.} Section~\ref{sec:qwen}'s
unpaired MTP drafter has been verified to pair correctly from GGUF metadata alone; it
has not been end-to-end tested serving actual multimodal-free inference, because the
build machine available had 7~GiB of RAM and could not hold both the text GGUF and the
mmproj simultaneously. That remains an open verification, not a known result.

None of the above is advanced as more than it is. This report is a record of three
models taken apart, the two bars that made ``unchanged'' a checked claim rather than an
assumed one, and one measurement that would not have been found by looking for it
directly --- it fell out of asking why a model that should obviously run faster
quantized sometimes does not.

\bibliography{references}
\bibliographystyle{plainnat}

\end{document}
